\section{Experiment 1 Supporting Statistics}
\label{app:supporting-multiplicity}

Table~\ref{tab:exp1-means} gives the absolute mode scores. The paired
differences in the main chapter remain the informative comparison.

\begin{table}[htbp]
    \centering
    \caption{Mean rubric score per mode, $n = 35$, with 95\,\% bootstrap
             confidence intervals ($B = 10\,000$).}
    \label{tab:exp1-means}
    \small
    \begin{tabular}{lcc}
        \toprule
        Mode & Total /25 [95\,\% CI] & KG-relevant /9 [95\,\% CI] \\
        \midrule
        \baseline{} & 8.03 [7.51, 8.54] & 4.66 [4.29, 5.03] \\
        \kg{}       & 8.74 [8.11, 9.37] & 5.34 [4.94, 5.74] \\
        \rag{}      & 9.06 [8.57, 9.54] & 5.09 [4.69, 5.49] \\
        \hybrid{}   & 8.60 [8.20, 9.00] & 5.11 [4.77, 5.46] \\
        \bottomrule
    \end{tabular}
\end{table}

Table~\ref{tab:exp1-deltas-full} gives the inferential details behind the
paired estimates in Table~\ref{tab:exp1-deltas}.

\begin{table}[htbp]
    \centering
    \caption{Paired permutation tests and effect sizes against \baseline{},
             $n = 35$ ($B = 20\,000$, two-sided, uncorrected).}
    \label{tab:exp1-deltas-full}
    \small
    \begin{tabular}{lcccccc}
        \toprule
        & \multicolumn{3}{c}{Total /25} & \multicolumn{3}{c}{KG-relevant /9} \\
        \cmidrule(lr){2-4} \cmidrule(lr){5-7}
        Mode & $\Delta$ [95\,\% CI] & $p$ & $d_z$ & $\Delta$ [95\,\% CI] & $p$ & $d_z$ \\
        \midrule
        \kg{}     & $+0.71$ [$+0.20$, $+1.26$] & 0.017 & $+0.44$ & $+0.69$ [$+0.29$, $+1.11$] & \textbf{0.003} & $+0.56$ \\
        \rag{}    & $+1.03$ [$+0.46$, $+1.60$] & 0.002 & $+0.59$ & $+0.43$ [$-0.03$, $+0.94$] & 0.127 & $+0.29$ \\
        \hybrid{} & $+0.57$ [$+0.09$, $+1.11$] & 0.056 & $+0.35$ & $+0.46$ [$+0.03$, $+0.91$] & 0.068 & $+0.34$ \\
        \bottomrule
    \end{tabular}
\end{table}

The governing Experiment~1 analysis is the pre-specified paired contrast
between each context arm and \baseline{}, tested by sign-flip permutation.
As a sensitivity to test choice and multiplicity, Table~\ref{tab:exp1-corrected}
also reports a Friedman omnibus and Bonferroni-corrected Wilcoxon signed-rank
tests. The Friedman test does not detect an omnibus mode difference on the
KG-relevant subscale. The corrected \kg{} comparison is therefore read only as
an alternative analysis of the planned contrast, not as post-hoc evidence
following a significant omnibus result.

\begin{table}[htbp]
    \centering
    \caption{Bonferroni-corrected Wilcoxon sensitivity. Friedman omnibus:
             $\chi^2 = 14.08$, $p = 0.003$ on total;
             $\chi^2 = 6.63$, $p = 0.084$ on KG-relevant.}
    \label{tab:exp1-corrected}
    \small
    \begin{tabular}{lcccc}
        \toprule
        & \multicolumn{2}{c}{Total /25} & \multicolumn{2}{c}{KG-relevant /9} \\
        \cmidrule(lr){2-3} \cmidrule(lr){4-5}
        Mode vs \baseline{} & $p$ raw & $p$ corrected & $p$ raw & $p$ corrected \\
        \midrule
        \kg{}     & 0.022 & 0.066 & 0.004 & \textbf{0.011} \\
        \rag{}    & 0.001 & \textbf{0.004} & 0.118 & 0.354 \\
        \hybrid{} & 0.048 & 0.143 & 0.064 & 0.192 \\
        \bottomrule
    \end{tabular}
\end{table}

\section{Experiment 1 Inter-Judge Agreement}
\label{app:supporting-agreement}

Table~\ref{tab:exp1-agreement} reports agreement before majority aggregation.
Pairs are evaluated on their shared valid rubric cells; Gemini abstained on
104 of 3\,500 cells, while the other judges returned a verdict for every cell.

\begin{table}[htbp]
    \centering
    \caption{Pairwise agreement in the final five-judge Experiment~1 panel.
             G4 = gpt-4o, GM = gemini-2.5-flash, CL = claude-sonnet-4.5,
             DS = deepseek-v4-pro, and GR = grok-4.6.}
    \label{tab:exp1-agreement}
    \small
    \begin{tabular}{lrrr}
        \toprule
        Pair & Valid cells & Raw agreement & Cohen's $\kappa$ \\
        \midrule
        G4--GM & 3\,396 & 85.7\,\% & 0.701 \\
        G4--CL & 3\,500 & 89.8\,\% & 0.788 \\
        G4--DS & 3\,500 & 90.3\,\% & 0.786 \\
        G4--GR & 3\,500 & 92.0\,\% & 0.822 \\
        GM--CL & 3\,396 & 85.8\,\% & 0.708 \\
        GM--DS & 3\,396 & 88.4\,\% & 0.753 \\
        GM--GR & 3\,396 & 87.3\,\% & 0.729 \\
        CL--DS & 3\,500 & 87.1\,\% & 0.727 \\
        CL--GR & 3\,500 & 87.5\,\% & 0.735 \\
        DS--GR & 3\,500 & 95.3\,\% & 0.893 \\
        \bottomrule
    \end{tabular}
\end{table}
% Source: results/JOERN_UNIFIED_FIVE_JUDGE.json, agreement section.

\section{Experiment 1 Robustness Statistics}
\label{app:supporting-robustness}

Table~\ref{tab:exp1-robustness} gives the statistics behind
Figure~\ref{fig:exp1-robustness}. The primary and single-judge rows use the
Joern CPG at $n = 35$. The generator and held-out rows are from an earlier
configuration that used a text-heuristic file matcher at $n = 40$ under a
three-judge panel; they are retained as directional evidence that alternative
generators preserve the sign of the effect but are not directly comparable
to the primary row. Values are therefore interpreted within rows.

\begin{table}[htbp]
    \centering
    \caption{The \kg{} effect under alternative generators, judges,
             and samples.}
    \label{tab:exp1-robustness}
    \small
    \setlength{\tabcolsep}{5pt}
    \begin{tabular}{llcccc}
        \toprule
        & & \multicolumn{2}{c}{Total /25} & \multicolumn{2}{c}{KG-relevant /9} \\
        \cmidrule(lr){3-4} \cmidrule(lr){5-6}
        Varied & Configuration & $\Delta$ & $p$ & $\Delta$ & $p$ \\
        \midrule
        ---       & primary configuration, $n = 35$ & $+0.71$ & \textbf{0.017} & $+0.69$ & \textbf{0.003} \\
        \midrule
        Judge     & gpt-4o alone             & --- & --- & $+0.63$ & \textbf{0.007} \\
                  & gemini-2.5-flash alone   & --- & --- & $+0.74$ & \textbf{0.005} \\
                  & claude-sonnet-4.5 alone  & --- & --- & $+0.26$ & 0.192 \\
                  & deepseek-v4-pro alone    & --- & --- & $+0.60$ & \textbf{0.019} \\
                  & grok-4.6 alone           & --- & --- & $+0.60$ & \textbf{0.009} \\
        \midrule
        Generator$^{*}$ & claude-haiku-4.5    & $+0.30$ & 0.407 & $+0.15$ & 0.457 \\
                  & gemini-2.5-flash         & $+0.33$ & 0.465 & $+0.47$ & 0.136 \\
                  & deepseek-chat            & $+0.47$ & 0.117 & $+0.20$ & 0.383 \\
        \midrule
        Sample$^{*}$ & held-out PRs, $n = 12$   & $+0.42$ & 0.426 & $+0.42$ & 0.312 \\
        \bottomrule
        \multicolumn{6}{l}{\footnotesize $^{*}$\,Text-heuristic builder, $n = 40$, three-judge panel; retained as directional evidence.} \\
    \end{tabular}
\end{table}

\section{Experiment 2 Detection Rates and Secondary Contrasts}
\label{app:supporting-exp2-contrasts}

Table~\ref{tab:exp2-detection} gives the exact detection rates behind
Figure~\ref{fig:exp2-detection}.

\begin{table}[htbp]
    \centering
    \caption{Detection rate on structural bands, $n = 28$, with Wilson
             95\,\% confidence intervals.}
    \label{tab:exp2-detection}
    \small
    \begin{tabular}{lcc}
        \toprule
        Arm & Detected & Rate [95\,\% CI] \\
        \midrule
        \baseline{}                 & 0/28  & 0.00 [0.00, 0.12] \\
        \rag{}                      & 5/28  & 0.18 [0.08, 0.36] \\
        \kg{} (deployed)            & 18/28 & 0.64 [0.46, 0.79] \\
        \hybrid{}                   & 21/28 & 0.75 [0.57, 0.87] \\
        \midrule
        \kg{} + inheritance edges   & 26/28 & 0.93 [0.77, 0.98] \\
        \kg{} idealised (ceiling)   & 26/28 & 0.93 [0.77, 0.98] \\
        \bottomrule
    \end{tabular}
\end{table}

Table~\ref{tab:exp2-perband} breaks detection down by structural band and
local control. The five structural bands (S1--S5) represent different
dependency types; the two local bands (L1--L2) serve as negative controls
whose consequences are visible in the diff.

\begin{table}[htbp]
    \centering
    \caption{Per-band detection counts. $n$ is the number of injections in
             each band. Structural bands (S1--S5) require cross-file
             reasoning; local bands (L1--L2) do not.}
    \label{tab:exp2-perband}
    \small
    \begin{tabular}{lrcccccc}
        \toprule
        Band & $n$ & \baseline{} & \kg{} & \rag{} & \hybrid{} & Inherit & Idealised \\
        \midrule
        S1 & 9  & 0/9  & 4/9  & 3/9  & 6/9  & 9/9  & 8/9  \\
        S2 & 9  & 0/9  & 6/9  & 1/9  & 7/9  & 7/9  & 8/9  \\
        S3 & 1  & 0/1  & 1/1  & 0/1  & 1/1  & 1/1  & 1/1  \\
        S4 & 6  & 0/6  & 4/6  & 1/6  & 4/6  & 6/6  & 6/6  \\
        S5 & 3  & 0/3  & 3/3  & 0/3  & 3/3  & 3/3  & 3/3  \\
        \midrule
        L1 & 6  & 6/6  & 6/6  & 6/6  & 6/6  & 6/6  & 6/6  \\
        L2 & 6  & 6/6  & 6/6  & 6/6  & 5/6  & 6/6  & 6/6  \\
        \bottomrule
    \end{tabular}
\end{table}

The structural bands confirm the aggregate pattern: baseline scores zero in
every band, and the eight cases recovered by inheritance edges fall in S1
($+5$), S2 ($+1$) and S4 ($+2$). The inheritance and idealised arms match
in aggregate, not band by band. Local controls are at or near ceiling for
all arms (one hybrid miss on L2).

The primary Experiment~2 contrast is \kg{} versus \baseline{}. Table
\ref{tab:exp2-secondary-contrasts} retains the secondary comparisons used to
interpret retrieval, hybrid concatenation, and inheritance edges.

\begin{table}[htbp]
    \centering
    \caption{Secondary exact McNemar contrasts on the 28 structural
             injections.}
    \label{tab:exp2-secondary-contrasts}
    \small
    \begin{tabular}{lccc}
        \toprule
        Contrast & $\Delta$ detection & Discordant wins/losses & $p$ \\
        \midrule
        \rag{} vs \baseline{} & $+0.179$ & 5/0 & 0.0625 \\
        \hybrid{} vs \kg{} & $+0.107$ & 4/1 & 0.3750 \\
        \kg{} + inheritance vs \kg{} & $+0.286$ & 8/0 & 0.0078 \\
        \bottomrule
    \end{tabular}
\end{table}

\subsection{Residual Idealised-Arm Misses}
\label{app:supporting-exp2-misses}

The two idealised-arm misses arise at different stages. In
\texttt{kafka\_S2\_01}, the review names genuine dependent classes, including
\texttt{CachingKeyValueStore} and \texttt{CachingSessionStore}, but does not
emit the exact file token required by the deterministic filename gate. This is
a boundary of the operational definition rather than an absence of affected
code in the review. In \texttt{kafka\_S1\_02}, whose resolver returns 82 true
dependents, the review warns generically about broken dependencies but names
none. That case is a generation or evidence-salience failure despite complete
structural evidence.
% Sources: experiments/2026-07-05_injection_exp2/out/manifest.json;
% out/reviews/{kafka_S2_01,kafka_S1_02}/kg_idealised.md;
% experiments/2026-09-23_five_judge_panel/RESULTS.json.

\section{Component Contributions}
\label{sec:results-ablations}
\label{app:supporting-components}

These component analyses support RQ2 by testing which parts of the graph block
contribute to cross-file detection. The rubric's regeneration variance
(Section~\ref{sec:repro}) makes it too coarse to separate sub-components, so
the analyses use the detection oracle instead.

The deployed graph arm supplies two kinds of dependency evidence
simultaneously: a file-level dependency list and resolved call-graph edges from
the code property graph. Table~\ref{tab:exp2-features} isolates each.

\begin{table}[htbp]
    \centering
    \caption{Detection on structural bands with one kind of dependency
        evidence at a time. Reviews use gpt-4o at $T=0.0$; rates are majority
        votes from five judges at $T=0.0$. The $p$~values are exact McNemar,
        Holm-corrected.}
    \label{tab:exp2-features}
    \begin{tabular}{llcc}
        \toprule
        Arm & Evidence supplied & Detected & Rate \\
        \midrule
        \baseline{}        & none                        & 0/28  & 0\,\% \\
        \texttt{deps only} & file-level dependency list   & 16/28 & 57\,\% \\
        \texttt{edges only}& resolved call edges         & 17/28 & 61\,\% \\
        \kg{}              & both                        & 18/28 & 64\,\% \\
        \midrule
        \multicolumn{2}{l}{Contrast} & $\Delta$ & $p$ (Holm) \\
        \midrule
        \multicolumn{2}{l}{\texttt{deps only} $-$ \baseline{}}  & $+57$pp & $< 0.001$ \\
        \multicolumn{2}{l}{\texttt{edges only} $-$ \baseline{}} & $+61$pp & $< 0.001$ \\
        \bottomrule
    \end{tabular}
\end{table}

Either component alone lifts detection from zero to over half, and both
survive correction. The dependency list recovers sixteen of the eighteen
combined-arm detections, while resolved call edges recover seventeen. Neither
component is significantly below the combined arm at $n=28$. The result
therefore supports partial substitutability, not the necessity or superiority
of resolved edges: the causal claim is that useful structural pointers enable
cross-file detection, while this sample cannot identify one representation as
uniquely responsible.

A pre-registered, Grafana/TypeScript-only test band, evaluated by the earlier
three-judge panel, addressed the test component. On 28 injections where a
renamed symbol provably breaks test files, the filename-convention test
matcher of Section~\ref{sec:approach-kg-render} named a broken test in 24 of 28 cases (86\,\%), compared with 27/28 (96\,\%)
when given the ground truth; the difference is not significant ($p = 0.38$).
Without the test
section, neither \baseline{} nor \kg{} named a broken test in any case. This
shows that the test block contributes detections the dependency block alone
does not.

\subsection{Context Relevance}
\label{sec:results-context-relevance}

The component ablation shows that dependency edges help, but does not rule
out a prompt-length confound: the reviewer might try harder whenever it sees
\emph{any} dependency block, regardless of content. A post-hoc control tests
this possibility.
Table~\ref{tab:exp2-signal-noise} tests this by holding the number of
context items constant and varying only their relevance.
% Source: experiments/2026-09-19_joern_replace_grep/context_ablation/RESULTS_SIGNAL_NOISE_V3.md
% Design follows Yang et al. (EMNLP 2025) GSM-DC matched-distractor methodology.

\begin{table}[htbp]
    \centering
    \caption{Detection under matched-quantity contexts that differ only in
        relevance. \texttt{noise} draws the same number of edges from
        unrelated symbols in the same repository; \texttt{relevant} filters
        the deployed edges to those pointing at true dependents.
        All arms use gpt-4o at $T=0.0$ and the earlier three-judge panel.}
    \label{tab:exp2-signal-noise}
    \begin{tabular}{llcc}
        \toprule
        Arm & Context relevance & Detected & Rate \\
        \midrule
        \baseline{}             & none             & 0/28  & 0\,\% \\
        \texttt{noise}          & wrong edges, matched count & 1/28  & 4\,\% \\
        \kg{}                   & deployed mix     & 18/28 & 64\,\% \\
        \texttt{relevant only}  & true-dependent edges only & 25/28 & 89\,\% \\
        \texttt{idealised}      & ground-truth resolver & 26/28 & 93\,\% \\
        \bottomrule
    \end{tabular}
\end{table}

The noise arm differs little from \baseline{} ($1/28$ versus $0/28$):
matched-volume edges drawn from unrelated symbols provide no comparable
detection gain. Filtering the deployed edges to only those pointing at true
dependents raises detection from 18 to 25 of 28. This post-hoc staircase is
consistent with the context-rot finding that
irrelevant tokens degrade output quality even when the answer is present in
the input~\cite{hong2025contextrot, gupta2025contextlength}.

\section{Human-Study Exploratory Endpoints}
\label{app:supporting-human}

The five non-primary criteria are exploratory and Holm-corrected. Only H5
(code clarity) remains distinguishable after correction, favouring the
\baseline{} review.

\begin{table}[htbp]
    \centering
    \caption{Exploratory human-study criteria. Values above $0.5$ favour
             \kg{}. IDs are prefixed H to distinguish them from the
             25-criterion rubric.}
    \label{tab:human-exploratory}
    \small
    \begin{tabular}{llcccc}
        \toprule
        ID & Criterion & Mean [95\,\% CI] & Non-tie $n$ & $p$ & $p_{\text{Holm}}$ \\
        \midrule
        H2 & Describes failure cases     & 0.469 [0.414, 0.528] & 25 & 0.146 & 0.499 \\
        H3 & Points to tests             & 0.454 [0.395, 0.512] & 19 & 0.125 & 0.499 \\
        H4 & Explains why                & 0.491 [0.438, 0.543] & 20 & 0.735 & 1.000 \\
        H5 & Comments on code clarity    & 0.414 [0.370, 0.460] & 24 & 0.006 & 0.028 \\
        H6 & Covers important issues     & 0.481 [0.417, 0.548] & 22 & 0.636 & 1.000 \\
        \bottomrule
    \end{tabular}
\end{table}

\section{Prompt and Review Length}
\label{app:supporting-length}

Table~\ref{tab:prompt-length} reports the total input size (system prompt
plus user prompt) and the augmentation overhead for each arm, computed
from all 40 stored evidence packs (including the five Go pull requests
excluded from the $n = 35$ analysis sample) and the generation template.
Table~\ref{tab:review-length} reports the output review length.

\begin{table}[htbp]
    \centering
    \caption{Total input size per arm ($n = 40$ stored evidence packs).
             Token estimates use the chars\,/\,4 heuristic.}
    \label{tab:prompt-length}
    \small
    \begin{tabular}{lrrrr}
        \toprule
        Arm & Mean chars & Median & $\approx$\,Mean tokens & Overhead vs \baseline{} \\
        \midrule
        \baseline{} & 16\,926 & 12\,338 & 4\,232 & --- \\
        \kg{}       & 18\,384 & 14\,240 & 4\,596 & $+$10.9\,\% \\
        \rag{}      & 21\,292 & 15\,733 & 5\,323 & $+$43.4\,\% \\
        \hybrid{}   & 22\,538 & 17\,016 & 5\,635 & $+$52.3\,\% \\
        \bottomrule
    \end{tabular}
\end{table}

\begin{table}[htbp]
    \centering
    \caption{Review output length per arm ($n = 40$ stored evidence packs).}
    \label{tab:review-length}
    \small
    \begin{tabular}{lrrr}
        \toprule
        Arm & Mean chars & Mean words & $\approx$\,Mean tokens \\
        \midrule
        \baseline{} & 1\,774 & 243 & 444 \\
        \kg{}       & 1\,991 & 266 & 498 \\
        \rag{}      & 1\,782 & 243 & 445 \\
        \hybrid{}   & 1\,857 & 252 & 464 \\
        \bottomrule
    \end{tabular}
\end{table}

The graph arm adds the least input overhead ($+$10.9\,\%) but produces
the longest reviews ($+$12\,\% in characters and $+$9\,\% in words relative
to baseline). RAG
adds more input ($+$43\,\%) yet produces reviews of the same length as
baseline. This asymmetry suggests the model converts structural context
into additional commentary more readily than it converts similarity
context.

